\section{Conclusion}

While APL and its descendant languages have attracted a devoted userbase,
there has been little cross talk between
the array-language community and
the broader programming language research community.
Much of the analysis opportunity taken for granted by lambda calculists
has been unavailable in APL---%
despite the many implementations,
such languages lacked formal semantics amenable to proofs.
Meanwhile, implementations of rank-polymorphic languages
have struggled with compilation due to control structure that is ``too dynamic,''
depending on run-time data to determine the structure of a loop nest.

Developing Remora's semantics kills two birds with one stone:
Formally stated reduction rules describe
the results expected from the implicit, data-driven control structure,
and typing rules give enough information about array shapes
to identify that control structure statically.
This necessarily entails recognizing programs
which cannot have such a control structure
due to incompatible array shapes---%
this is not what we set out to do
but a benefit realized by pursuing a larger goal.
By casting the aggregate lifting as an extension to
$\lambda$-calculus's function application semantics,
we escape from APL's limitation on function arity
and can treat functions as first-class values.

There are two things to look for in evaluating a type system.
We want to know the conclusions drawn by type checking reflect reality,
which is handled by a conventional type soundness theorem.
We also want to be sure that the types convey the information we seek.
In Remora's case, that information is
the iteration structure implicit in each function application.
The typing rule for function application
(and similarly, the rules for type and index application)
produces a static characterization of that implicit iteration structure.

Our type erasure can be seen as a compilation pass
which moves the decision about how to break arguments into cells
from the function's type into the application term.
While our primary purpose in presenting type erasure
is to point out type-level information which is not truly needed at run time,
it also serves to make the control structure one step more explicit.
Arguments' full shapes---available both at erasure time
and later by inspecting argument terms---%
suffice to determine the arguments' frames,
the last puzzle piece needed to turn
Remora's implicit iteration structure into
explicit calls to {\tt map} and {\tt replicate} functions.
Our intention for future work is to
demonstrate use of that shape information for
fully static compilation of Remora code.